% !TeX root = ../main.tex

\section{Solving Sudokus With Brainscales-2}

We will apply the simulated neurons to solve a given Sudoku problem. For an \( n \times n \) Sudoku grid, there are \( n^2 \) fields which could each be \( n \) different numbers. We will use \( n \) neurons for each field. These will each represent one of the possible output numbers, using a \quotedef{one hot encoding}.

We will use the Brainscales-2 system within Python to set up the neurons and synapses in order to solve the Sudoku. The created and used code is included as attached \texttt{.ipynb} and \texttt{.py} files.

\subsection{Implementation}

We set up neurons to represent the board state. We also use additional neurons to add random noise to the calculation (this avoids dead-ends where a mistake in the solve happened). These are connected by synapses.

In order to sustain activity and \quoteapprox{save} a value when a number is determined, neurons are linked to themselves.

The constraints/rules of the Sudoku are encoded in other synapses between different groups of neurons.

These will be explained for a \( 4 \times 4 \) and a \( 9 \times 9 \) Sudoku grid. The exact number of neurons and synapses needed will be calculated as well.

\subsubsection{Summary of Neurons and Synapses for a \( 4 \times 4 \) Sudoku Grid}

For a \( 4 \times 4 \) Sudoku, 64 neurons are needed to represent the state of the board.

In addition to this, 64 stimulating input neurons/spikes are used.
Each of these is connected to one neuron, meaning 64 stimulating synapses.

There is 1 additional stimulating input spiking neuron connected to all neurons.
The synapse weights for this neuron are adjusted to 0 or 63 depending on if the corresponding clue was given or not.
This leads to a maximum of \( 4^2 = 16 \) \quoteapprox{active} synapses with weight \( ≠ 0 \), if each field has a clue.

Constraint Synapses:
\begin{itemize}
    \item The constraint for one number per field means for each of the \( 4^2 = 16 \) fields, inhibitory synapses are needed to connect each of the 4 neurons representing each number to the other 3.
    \subitem This means \( 4^2 \cdot 4 \cdot 3 = 192 \) synapses.

    \item The constraint for no repeat numbers in each row means for each of the \( [4 \ \text{numbers} \cdot 4 \ \text{rows} = 16] \) neurons in rows,
    inhibitory synapses are needed connecting each of the 4 neurons to the other 3.
    \subitem This means \( 4^2 \cdot 4 \cdot 3 = 192 \) synapses.

    \item The constraint for no repeat numbers in each column similarly means each of the \( [4 \ \text{numbers} \cdot 4 \ \text{rows} = 16] \) neurons columns,
    inhibitory synapses connect each of the 4 neurons to the other 3.
    \subitem This means \( 4^2 \cdot 4 \cdot 3 = 192 \) synapses.

    \item The last constraint calls for no repeat numbers in each of the \( [2^2 = 4] \) blocks with shape: \( (2 x 2) \), as the rules of Sudoku dictate.
    This constraint applies to each number, so there are in total \( [2^3 = 8] \) blocks of 4 neurons.
    Inhibitory synapses must connect each of the 4 neurons in each block (with same number) to the other 3.
    \subitem This means \( 2^3 \cdot 4 \cdot 3 = 96 \) synapses.
\end{itemize}

Concluding the summary:

There are 64 main neurons doing the calculations.
Counting the \( [64 + 1 = 65] \) stimulating input spike neurons leads to a total of:
\[129 \ \text{neurons.}\]
The calculations considering \( [64 + 16 = 80] \) stimulating neurons and 672 inhibiting neurons lead to a total of:
\[752 \ \text{synapses.}\]


\subsubsection{Summary of Neurons and Synapses for a \( 9 \times 9 \) Sudoku Grid}

For a \( (9 \times 9) \) Sudoku, 729 neurons are needed to represent the state of the board.

In addition to this, 729 stimulating input neurons/spikes are used.
Each of these is connected to one neuron, meaning 729 stimulating synapses.

There is 1 additional stimulating input spiking neuron connected to all neurons.
The synapse weights for this neuron are adjusted to 0 or 63 depending on if the corresponding clue was given or not.
This leads to a maximum of \( 9^2 = 81 \) "active" synapses with weight \( ≠ 0 \), if each field has a clue.

Constraint Synapses:
\begin{itemize}
    \item The constraint for one number per field means for each of the \( 9^2 = 81 \) fields,
    inhibitory synapses were needed to connect each of the 9 neurons representing each number to the other 8.
    \subitem This means \( 9^2 \cdot 9 \cdot 8 = 5832 \) synapses.

    \item The constraint for no repeat numbers in each row means for each of the \( [9 numbers \cdot 9 rows = 81] \) neurons in rows,
    inhibitory synapses are needed connecting each of the 9 neurons to the other 8.
    \subitem This means \( 9^2 \cdot 9 \cdot 8 = 5832 \) synapses.

    \item The constraint for no repeat numbers in each column similarly means each of the \( [9 numbers \cdot 9 rows = 81] \) neurons columns,
    inhibitory synapses connect each of the 9 neurons to the other 8.
    \subitem This means \( 9^2 \cdot 9 \cdot 8 = 5832 \) synapses.

    \item The last constraint calls for no repeat numbers in each of the \( [3^2 = 9] \) blocks of shape \( (3 \times 3) \), as the rules of Sudoku dictate.
    This constraint applies to each number, so there are in total \( [3^3 = 27] \) blocks of 9 neurons.
    Inhibitory synapses must connect each of the 9 neurons in each block (with same number) to the other 8.
    \subitem This means \( 3^3 \cdot 9 \cdot 8 = 1944 \) synapses.
\end{itemize}

Concluding the summary:

There are 729 main neurons doing the calculations.
Counting the \( [729 + 1 = 730] \) stimulating input spike neurons leads to a total of:
\[1459 \ \text{neurons.}\]
The calculations considering \( [729 + 81 = 810] \) stimulating neurons and 19440 inhibiting neurons lead to a total of:
\[20250 \ \text{synapses.}\]